\section{A Genealogy of the Criterion}

Positions are examined; movements are also \emph{lived}, and the living leaves a record. The exclusivist thesis is unusual among philosophical doctrines in that its career is itself part of the evidence about it: it was announced as a law of history, organized as a movement, engineered into a formal criterion, repaired under fire, and finally abandoned by its own heirs---all within about a century, and almost all of it in documents that can still be read. This section reads them. Nothing in it yet argues that the thesis is false; the argument comes after. What the genealogy establishes is something an examination needs first: what the position actually claimed, in its own strongest words, and what happened when its own adherents tried to make the claim precise.

\subsection{Comte: The Law That Explained Its Rivals Away}

Comte's opening move, quoted in the first section, was to announce a ``great fundamental law'' of intellectual history. It is worth watching how much work the announcement does. The three states are not presented as three schools among which one must argue; they are presented as ages. ``The first is the necessary point of departure of the human understanding; and the third is its fixed and definitive state. The second is merely a state of transition.''\endnote{Comte, in Martineau's condensation, \emph{The Positive Philosophy of Auguste Comte}, vol.~1 (London: George Bell \& Sons, 1896), Introduction, ch.~I, p.~2. Wording checked 2026-07-25 in the same Internet Archive scan of the 1896 Bell printing used throughout (public domain). Obvious OCR artifacts in the scan are silently normalized to the printed reading only where the reading is unambiguous.} In the theological state the mind ``supposes all phenomena to be produced by the immediate action of supernatural beings''; in the metaphysical, by ``abstract forces, veritable entities (that is, personified abstractions)''; in the positive state it studies only laws of phenomena, ``their invariable relations of succession and resemblance.'' The proof offered for the law is partly historical---every science, Comte says, shows marks of having passed through the earlier stages---and partly a charming appeal to biography: ``each of us is aware, if he looks back upon his own history, that he was a theologian in his childhood, a metaphysician in his youth, and a natural philosopher in his manhood.''\endnote{Martineau, vol.~1, pp.~2--3, checked 2026-07-25 as above. The biographical argument is offered by Comte as ``an indirect evidence'' of the general law. Its force may be measured by trying it in reverse: an age that reasons from its own maturity to the falsity of its childhood convictions has assumed, not shown, that the development was an education rather than a forgetting.}

Two features of the system deserve more attention than the slogan usually receives. The first is a concession buried in Comte's own argument for the law. Explaining why the mind \emph{had} to begin theologically, he writes: ``If it is true that every theory must be based upon observed facts, it is equally true that facts cannot be observed without the guidance of some theory. Without such guidance, our facts would be desultory and fruitless; we could not retain them: for the most part we could not even perceive them.''\endnote{Martineau, vol.~1, p.~4, checked 2026-07-25. Comte's point is genetic---primitive man needed theology as scaffolding to escape ``a vicious circle'' between theory and observation---but the epistemological premise he states is general, and it is the premise this essay's third section will need: observation is never a self-sufficient tribunal, because what counts as a fact is already the work of a theorizing mind.} Comte deploys the point to explain the past; he does not notice that it undermines the future he predicts, in which observation alone would sit as judge. If facts cannot even be \emph{perceived} without theory, then the positive stage never possesses the theory-free ground the schema promises, and the renunciation of everything beyond ``invariable relations'' is itself one more theoretical stance, adopted and defensible---if at all---only by philosophical argument of exactly the kind the schema declares obsolete.

The second feature is the system's ambition, which was never merely descriptive. The positive philosophy, Comte announces, ``offers the only solid basis for that Social Reorganization which must succeed the critical condition in which the most civilized nations are now living''; the crisis of the age ``is shown by a rigid analysis to arise out of intellectual anarchy,'' and order will return when one of the three warring philosophies prevails---the positive, since it alone has been advancing.\endnote{Martineau, vol.~1, pp.~15--17, checked 2026-07-25. The chapter is explicit that ``all social mechanism rests upon Opinions'' and that the disorder of the age is the coexistence of ``three incompatible philosophies''; stability requires the triumph of one. The task of ``reorganizing society'' is said to be ``too mighty'' for either theology or metaphysics.} This is worth remembering whenever positivism presents itself as modesty about what can be known. The founding document is not modest. It proposes the government of opinion by a single doctrine, promises that agreement on first principles will end revolutionary disorder, and assigns to a philosophy of science the office that theologies once held. The criterion of the knowable was born as an instrument of rule.

And it made testable predictions. Applying his method to astronomy, Comte drew the boundary of the knowable with complete confidence: of the planets ``we can never know anything of their chemical or mineralogical structure; and, much less, that of organized beings living on their surface''; and as to the temperatures of the heavenly bodies, ``I regard this order of facts as for ever excluded from our recognition.''\endnote{Martineau, vol.~1, Book~II (Astronomy), ch.~I, pp.~148--149, checked 2026-07-25 in the same scan. Comte's stated rule was that whatever knowledge depends on senses other than sight ``we must at once exclude from our expectations''; stellar chemistry and temperature fell, he judged, on the far side of the line.} The sequel is among the best-known reversals in the history of science: within a generation, the spectral analysis founded by Kirchhoff and Bunsen (1859--1861) was reading the chemical composition of the sun and stars in their light, and stellar temperatures followed.\endnote{Reported as the standard history of spectroscopy: Kirchhoff's identification of solar absorption lines with terrestrial elements (1859) and the Kirchhoff--Bunsen spectral analyses of 1859--61 founded solar and stellar chemistry. No dedicated witness was collated for this common historical fact; it is load-bearing here only as an illustration, and the illustration's force survives any refinement of the dates.} The moral is not that Comte was careless; he was applying his criterion exactly. The moral is that a priori boundaries of the knowable, drawn with the full authority of a scientific philosophy, were falsified by science itself---and falsified, ironically, through Comte's own favored sense of sight, by an inference from visible spectral lines to invisible constitution. Reality outran the rule; the instruments themselves refuted the philosophy of the instruments. A criterion that had misjudged what observation could reach within the sciences might reasonably have been suspected of misjudging what reason could reach beyond them. It was not; it was sharpened instead.

\subsection{Vienna: A Movement with a Manifesto}

The sharpening had a place, a date, and a pamphlet. In 1929 the Verein Ernst Mach in Vienna published a fifty-nine-page programmatic tract, \emph{Wissenschaftliche Weltauffassung: Der Wiener Kreis} (``The Scientific Conception of the World: The Vienna Circle''), naming no author on its title page; its foreword is signed, for the society, by Hans Hahn, Otto Neurath, and Rudolf Carnap, and it was presented to Moritz Schlick, around whose seminar the Circle had formed.\endnote{\emph{Wissenschaftliche Weltauffassung: Der Wiener Kreis} (Vienna: Artur Wolf Verlag, 1929; Ver\"offentlichungen des Vereines Ernst Mach). Publication facts and all quoted wording checked 2026-07-25 at the transcription maintained by the Collected Works of Rudolf Carnap editorial project, phil.cmu.edu (the transcription notes ``Das Werk hat keine offiziell genannten Verfasser''---the work has no officially named authors---and prints the foreword over the names Hahn, Neurath, Carnap). The German text, first published in 1929, is in the public domain in the United States; translations of it here are working glosses by this essay, and the German governs. The standard scholarly attribution of the drafting to Neurath with Hahn and Carnap is as reported in Thomas Uebel, ``The Vienna Circle,'' \emph{Stanford Encyclopedia of Philosophy} (substantive revision 7 May 2024), checked 2026-07-25.} The document is brief, confident, and remarkably candid about what the movement was and wanted. It is the nearest thing the exclusivist thesis possesses to a creed, and its sentences deserve exact quotation.

On the texture of reality: ``In der Wissenschaft gibt es keine `Tiefen'; \"uberall ist Oberfl\"ache''---in science there are no ``depths''; everywhere is surface. On mystery: ``Die wissenschaftliche Weltauffassung kennt keine unl\"osbaren R\"atsel''---the scientific world-conception knows no insoluble riddles; traditional philosophical problems are partly unmasked as pseudo-problems, partly converted into empirical questions. On method against the theist: when someone asserts ``there is no God,'' or that the ground of the world is the unconscious, ``we do not say to him: `what you say is false'; rather, we ask him: `what do you mean by your statements?'\,''---and the interrogation is expected to sort his sentences into the empirically redeemable and the ``v\"ollig bedeutungsleer,'' the completely empty of meaning.\endnote{Manifesto, part~II (``Die wissenschaftliche Weltauffassung''), checked 2026-07-25 as above; glosses by this essay. Note the exact anticipation of Ayer's tactic seven years later: the position's force lies wholly in replacing the question of truth by the question of meaning, and its examples of condemned sentences include, side by side, the atheist's denial and the theist's assertion.} The metaphysician and the theologian, the manifesto explains, ``glauben, sich selbst mi{\ss}verstehend, mit ihren S\"atzen etwas auszusagen, einen Sachverhalt darzustellen''---believe, misunderstanding themselves, that their sentences state something, represent a state of affairs---whereas analysis shows these sentences express only a \emph{Lebensgef\"uhl}, an attitude toward life, whose adequate medium is not theory but art, ``zum Beispiel Lyrik oder Musik.'' What such men produce is ``nicht Theorie\ldots\ sondern Dichtung oder Mythus''---not theory, but poetry or myth, mislabeled.

The manifesto is equally frank about its program. There is no philosophy as a fundamental or universal science above the empirical sciences; ``es gibt keinen Weg zu inhaltlicher Erkenntnis neben dem der Erfahrung; es gibt kein Reich der Ideen, das \"uber oder jenseits der Erfahrung st\"ande''---no path to substantive knowledge besides experience, no realm of ideas above or beyond it. The workers of the scientific world-conception set themselves ``mit Vertrauen''---with confidence---to the task of clearing away ``den metaphysischen und theologischen Schutt der Jahrtausende,'' the metaphysical and theological rubble of millennia. And the closing line joins the epistemology to a social hope in a single breath: ``Die wissenschaftliche Weltauffassung dient dem Leben und das Leben nimmt sie auf''---the scientific world-conception serves life, and life receives it.\endnote{Manifesto, part~IV (``R\"uckblick und Ausblick'') and closing sentence, checked 2026-07-25; glosses by this essay. The same closing pages read the growth of metaphysical and theological ``inclinations'' sociologically, as an artifact of class struggle and disappointed masses, and predict that a production process ever more shaped by the machine ``leaves ever less room for metaphysical ideas''---a prediction whose fortunes the reader may assess independently.} Comte's structure is intact beneath the modern logic: a boundary of the meaningful policed in the name of science, a philosophy of history in which the excluded discourse is debris left by an earlier age, and an explicit social mission. What Vienna added was the instrument the nineteenth century lacked: the new mathematical logic of Frege and Russell, which seemed to promise that the boundary could at last be drawn \emph{formally}---not as a sage's verdict but as a theorem-shaped rule anyone could apply. That promise is what the next two decades tested.

It should be said, because the essay will argue hard against these men's doctrine: the Circle's members were serious scientists and logicians, and their fate was not comfortable. Hahn died in 1934; Schlick was murdered on the university steps in 1936; the Verein Ernst Mach was dissolved for political reasons; within a few years emigration and exile had scattered the rest.\endnote{Reported from Uebel, ``The Vienna Circle,'' \emph{Stanford Encyclopedia of Philosophy} (rev.\ 2024), checked 2026-07-25, which tabulates the deaths, emigrations, and the 1936 murder of Schlick. The dispersal belongs to the history of European catastrophe, not to the argument of this essay; it is recorded here because the movement's afterlife---logical empiricism as the dominant philosophy of science of the mid-century English-speaking world---was conducted by \'emigr\'es, and because candor about an opponent's history is part of taking the opponent seriously.} The doctrine they carried into English-speaking philosophy was examined there with full rigor---by friends. That examination is the next part of the story.

\subsection{Ayer: The Criterion Engineered}

\emph{Language, Truth and Logic} is the criterion's engineering document, and its opening page announces the ambition without preliminaries: ``The traditional disputes of philosophers are, for the most part, as unwarranted as they are unfruitful.''\endnote{Ayer, \emph{Language, Truth and Logic}, ch.~1, opening sentence; checked 2026-07-25 in the same Internet Archive scan of the Pelican reprint identified earlier. In the preface Ayer locates himself exactly: his views ``derive from the doctrines of Bertrand Russell and Wittgenstein, which are themselves the logical outcome of the empiricism of Berkeley and David Hume,'' and the philosophers with whom he is ``in the closest agreement'' are ``the `Viennese circle'\,''; among them he owes most to Carnap. The lineage from the flames passage through Vienna to the criterion is thus not this essay's construction but Ayer's own genealogy of himself.} The preface divides all genuine propositions, ``like Hume,'' into relations of ideas and matters of fact, and states the test: an empirical hypothesis need not be conclusively verifiable, but ``some possible sense-experience should be relevant to the determination of its truth or falsehood.'' A putative proposition failing this and not being a tautology ``is metaphysical, and\ldots\ being metaphysical, it is neither true nor false but literally senseless.''

The engineering interest of the book lies in what happens when Ayer tries to give the rule a precise form, because every load-bearing decision was forced on him by the sciences he meant to protect. Conclusive verifiability---the ``strong'' sense, in which a proposition is verifiable ``if, and only if, its truth could be conclusively established in experience''---had to be rejected: universal laws of nature (``arsenic is poisonous''; ``all men are mortal'') cover infinitely many cases and can never be conclusively established, so a strong criterion ``will prove too much,'' condemning the whole of science and history along with theology; indeed it is ``self-stultifying as a criterion of significance,'' since on Ayer's own view no proposition beyond a tautology is more than a probable hypothesis.\endnote{Ayer, \emph{Language, Truth and Logic}, ch.~1, the strong/weak discussion; wording checked 2026-07-25 in the same scan. Ayer notes that some positivists (he cites Schlick and Waismann) had accepted the ``heroic course'' of calling laws of nature nonsense---``albeit an essentially important type of nonsense''---and replies dryly that ``important'' is ``simply an attempt to hedge.''} Conclusive \emph{falsifiability} fared no better: a hypothesis can always be saved from refuting observations by blaming auxiliary assumptions, so definite confutation is as unattainable as definite proof---and Ayer's footnote names the proposal he is rejecting as that of ``Karl Popper in his \emph{Logik der Forschung}.''\endnote{Ayer, \emph{Language, Truth and Logic}, ch.~1, with the footnote citing Popper's \emph{Logik der Forschung}; checked 2026-07-25 in the same scan. The passage matters doubly: it documents that the falsifiability idea was on the table from the beginning, and it documents why the positivists could not use it for their purpose---which was a boundary of \emph{meaning}, not of science. What Popper actually proposed is the subject of the next subsection.} So Ayer settled on the weak criterion: a statement is significant if, in conjunction with ``certain other premises,'' experiential propositions can be deduced from it that do not follow from those premises alone. ``This criterion seems liberal enough,'' he remarks. It was---fatally so, as he would be the one to admit.

\subsection{Popper: The Alternative That Refused to Be an Ally}

Karl Popper is often conscripted into this story as one more verificationist, and the conscription reverses his actual position, which he spent decades protesting. His problem, formed in the Vienna of the 1920s though he was never a member of the Circle, was not the positivists' problem. They sought a criterion of \emph{meaning}: a line between sense and nonsense. He sought a criterion of \emph{demarcation}: a line between empirical science and everything else---a line, as he put it, between theories like Einstein's, which forbid observable outcomes and so run risks, and theories that can absorb any observation whatever. From that contrast he drew his famous conclusions: ``Every `good' scientific theory is a prohibition: it forbids certain things to happen''; and ``A theory which is not refutable by any conceivable event is nonscientific. Irrefutability is not a virtue of a theory (as people often think) but a vice.''\endnote{Karl Popper, ``Science: Conjectures and Refutations,'' the 1953 lecture printed as ch.~1 of \emph{Conjectures and Refutations: The Growth of Scientific Knowledge} (London: Routledge \& Kegan Paul, 1963). The work is in copyright; brief quotations appear here under ordinary fair-quotation practice. Wording was checked 2026-07-25 at an unofficial web transcription of the chapter; the transcription is a wording-check witness only, no printed edition was collated for this essay, and the printed text governs. Popper's \emph{Logik der Forschung} (Vienna: Springer, imprinted 1935, issued 1934), translated as \emph{The Logic of Scientific Discovery} (London: Hutchinson, 1959), is cited bibliographically, not collated; the account of its doctrine follows the chapter quoted here and Stephen Thornton, ``Karl Popper,'' \emph{Stanford Encyclopedia of Philosophy} (substantive revision 12 September 2022), checked 2026-07-25.}

The distinction the instrumentalist needs to blur, Popper states with every resource of emphasis. Of the boundary he drew he says: it ``was neither a problem of meaningfulness or significance, nor a problem of truth or acceptability. It was the problem of drawing a line\ldots\ between the statements, or systems of statements, of the empirical sciences, and all other statements.'' Of the meaning question itself: ``I personally was never interested in the so-called problem of meaning; on the contrary, it appeared to me a verbal problem, a typical pseudo-problem.'' Of the misreading: ``Neither falsifiability nor testability were proposed by me as criteria of meaning.'' And---decisively for this essay's subject---of what falls outside science: ``if a theory is found to be non-scientific, or `metaphysical'\ldots\ it is not thereby found to be unimportant, or insignificant, or `meaningless', or `nonsensical'.'' Myths, he adds, may develop into testable science; ``historically speaking all---or very nearly all---scientific theories originate from myths, and\ldots\ a myth may contain important anticipations of scientific theories.''\endnote{Popper, ``Science: Conjectures and Refutations,'' as in the preceding note. Popper's examples of fruitful metaphysics include Empedocles' evolutionary conjecture and the atomism that preceded testable atomic theory; his examples of pseudo-science protected by irrefutability were, notoriously, Marxist historical theory and the psychoanalysis of Freud and Adler. Nothing in this essay depends on endorsing those particular verdicts, or on the adequacy of falsifiability itself as a demarcation---a question with its own critical literature (the roles of auxiliary hypotheses and of degrees of testability among the classic difficulties, as reported in Thornton's SEP article). What the essay takes from Popper is only what he insists on: that the line around science is not a line around sense.}

Note what this does to the instrumentalist's rhetorical economy. ``Falsifiability'' is the piece of philosophy of science most likely to be quoted at a believer, usually as though it were a hygiene rule descended from verificationism. In its author's hands it was the opposite of the verificationists' rule in the respect that matters most here: it demarcates without condemning. Popper judged the verifiability criterion a failure even as demarcation---``it excludes from science practically everything that is, in fact, characteristic of it (while failing in effect to exclude astrology)''---and he refused, in terms, to let the boundary of science double as a boundary of the meaningful or the knowable. The instrumentalist who retreats from Ayer to Popper has therefore not found a fallback position; he has surrendered the disputed territory. Under Popper's flag, the question of God may be non-science. It is not thereby nonsense, and the volume must still be read before anyone presumes to burn it.

\subsection{The Retreat of 1946}

Ten years after the first edition, Ayer reissued \emph{Language, Truth and Logic} with a long new introduction that is one of the most instructive documents in twentieth-century philosophy---a fortification commander's tour of his own breached walls, conducted with complete honesty. The opening concession has been quoted already: the book's questions ``are not in all respects so simple as it makes them appear.'' The details are what matter.

First, the strong/weak distinction of 1936 dissolves under his own hand: as he had represented them, he now admits, ``these are not two genuine alternatives,'' since on his own doctrine no proposition beyond a tautology is ever conclusively established, so the ``strong'' sense had no possible application.\endnote{Ayer, introduction to the 2nd edition (1946), section ``The Principle of Verification''; wording checked 2026-07-25 in the same scan. He partially rescues conclusive verification for ``basic propositions'' recording a single present experience, then observes that such propositions convey nothing beyond the experience itself: ``It is, in short, a case of `nothing venture, nothing lose'\ldots\ equally a case of `nothing venture, nothing win'.''} Second---and this is the sentence on which the technical collapse turns---he concedes that the weak criterion of 1936, the one advertised as ``liberal enough,'' in fact ``allows meaning to any statement whatsoever.'' His own counterexample: given any nonsense sentence \emph{S}---his choice was ``the Absolute is lazy''---and the auxiliary premise ``if \emph{S} then this is white,'' the observation-statement ``this is white'' follows from the pair and not from the auxiliary alone; so \emph{S} passes the test, ``and this would hold good for any other piece of nonsense that one cared to put\ldots\ in its place.'' The footnote credits the objection to Isaiah Berlin.\endnote{Ayer, 1946 introduction, same section, checked 2026-07-25; the footnote cites Berlin's ``Verifiability in Principle,'' \emph{Proceedings of the Aristotelian Society} 39 (1938--39). Ayer immediately adds that ``the same objection applies to the proposal that we should take the possibility of falsification as our criterion,'' since any statement conjoined with a suitable conditional contradicts some observation-statement---the mirror-image failure. A criterion of meaning, unlike Popper's criterion of demarcation, cannot tolerate these results, because it was supposed to sort \emph{all} discourse.} Ayer then rebuilt the criterion a second time, in the two-story form of ``directly'' and ``indirectly'' verifiable statements, with clauses restricting the auxiliary premises. Three years later Alonzo Church's review proved, in half a page of the \emph{Journal of Symbolic Logic}, that the rebuilt criterion again admitted (near enough) any statement whatever.\endnote{Alonzo Church, review of \emph{Language, Truth and Logic}, 2nd ed., \emph{Journal of Symbolic Logic} 14 (1949): 52--53; cited bibliographically, not collated; the standard account of its result is as reported in the Stanford Encyclopedia narratives already identified. The pattern---repair, breach, repair, breach---is the documented career of every sharp formulation of the criterion.}

By mid-century the salvage operation had an auditor. Carl Hempel's 1950 survey, identified in the previous section, walked the whole sequence---conclusive verifiability, conclusive falsifiability, the weak criteria, translatability into an empiricist language---and reported that each either condemned manifest science or condemned nothing; his conclusion was that cognitive significance admits of degrees and cannot be captured by any sharp criterion of the promised kind.\endnote{Hempel, ``Problems and Changes in the Empiricist Criterion of Meaning'' (1950), with the 1965 revision, as identified bibliographically in the earlier section. The precise verification ceiling asserted here should be stated exactly: Hempel's texts were identified but not read for this essay; no Hempel wording is quoted anywhere in it; the account of his survey's structure (which criteria he examined, in which order, and the degrees-of-significance conclusion) is taken from the named secondary witness, James Fetzer's ``Carl Hempel,'' \emph{Stanford Encyclopedia of Philosophy} (substantive revision 16 March 2026), checked 2026-07-24 and 2026-07-25, corroborated by Ayer's own 1946 admissions quoted above, which were read directly. Any claim about Hempel finer than this---his examples, his exact formulations, his qualifications---is deliberately not made.} And Ayer's own final position, stated in the 1946 introduction and examined at length in the next section, was that the principle should be regarded ``not as an empirical hypothesis, but as a definition,'' though ``not\ldots\ entirely arbitrary''---an offer whose costs will be counted where it is examined. Understatement is also a philosophical genre.

\subsection{Quine: The End of Sentence-by-Sentence Meaning}

The last blow came from inside the empiricist family and was aimed at the foundations rather than the formulations. W.~V. Quine's ``Two Dogmas of Empiricism'' (1951) attacked the two pillars on which any Ayer-style criterion must stand: the analytic/synthetic distinction---the neat split between truths of meaning and truths of fact that gave the criterion its first disjunct---and ``reductionism,'' the belief that each meaningful statement, taken singly, has its own fund of confirming experiences. Against the second, Quine set his celebrated holism: our statements ``face the tribunal of sense experience not individually but only as a corporate body''; ``the unit of empirical significance is the whole of science.'' The totality of our knowledge, ``from the most casual matters of geography and history to the profoundest laws of atomic physics or even of pure mathematics and logic,'' is ``a man-made fabric which impinges on experience only along the edges.'' It follows that ``any statement can be held true come what may, if we make drastic enough adjustments elsewhere in the system,'' and, ``by the same token, no statement is immune to revision.''\endnote{W.~V. Quine, ``Two Dogmas of Empiricism,'' \emph{Philosophical Review} 60 (1951): 20--43, revised in \emph{From a Logical Point of View} (1953). The paper is in copyright; it is quoted here minimally under ordinary fair-quotation practice. Wording was checked 2026-07-25 at the widely used web transcription edited by Andrew Chrucky (ditext.com), which marks the 1951 and 1961 readings; the transcription is a wording-check witness, and the printed journal text governs. The previous edition of this essay cited the paper from its standard bibliographic record only; the deepened edition upgrades the citation to checked quotation at the named transcription, and says so.}

For the criterion, the consequence is terminal in a way the earlier objections were not. Berlin and Church showed that particular formulations leaked; Quine argued that there is nothing for a correct formulation to \emph{be}, because individual sentences---save perhaps those at the observational periphery---do not have their own separable empirical content to test. If confirmation attaches to the web and not the thread, a thread-by-thread meaning test is not hard to state but wrong in kind. The philosophical world largely agreed, and the verification principle passed out of live doctrine---not refuted by its enemies, but dismantled by Berlin, Church, and Quine and pronounced upon by Hempel, empiricists all. What Quine's web costs the \emph{instrumentalist}, as distinct from the criterion, is examined where the essay confronts the holist reply directly, in the section on the laboratory; the short version is that a web in which metaphysics, mathematics, and physics differ in degree of entrenchment rather than in kind of meaning is a web in which the metaphysician is again a colleague to be argued with, not a patient to be diagnosed.

The genealogy is now complete, and one feature of it should be fixed before the analysis begins, because the analysis will explain it. Every fatal wound was self-inflicted, or inflicted by kin: Comte's boundary of the knowable was broken by the sciences it claimed to describe; Ayer condemned the strong criterion himself; Berlin, Church, and Quine were empiricists; Hempel was a logical empiricist writing the movement's audit from within; Popper, the outsider nearest at hand, rejected not the criterion's answer but its question. The historical record shows \emph{that} the criterion could not be made to work. It does not yet show \emph{why} no repair was ever going to succeed. For that, the criterion must be made to face the one sentence it was never designed to sort: itself.
